\Lecture{Regular expressions}
Regular expressions are closely related to NFA's. 

\begin{definition}[Regular expessions]. \label{def:re}

    \textbf{Basis:} 
    \begin{itemize}
        \item The constants $\epsilon$ and $\emptyset$ are regular expressions, for the languages $L(\epsilon) = \set{\epsilon}$ and $L(\emptyset) = \emptyset$.
        \item If $a$ is any symbol, then $a$ is a regular expression for the language $L(a) = \set{a}$.
        \item A variable such as $L$, is a variable representing any language. 
    \end{itemize}
    \textbf{Induction}
    \begin{itemize}
        \item If $E$ and $F$ are regular expressions, then $E + F$ is a regular expression for the language $L(E + F) = L(E) \cup L(F)$.
        \item If $E$ and $F$ are regular expressions, then $EF$ is a regular expression for the language $L(EF) = L(E)L(F)$ (concatenation). 
        \item If $E$ is a regular expression, then $E^*$ is a regular expression for the language $L(E^*) = (L(E))^*$.
        \item If $E$ is a regular expression, then $(E)$ is also a regular expression for the same language $L((E)) = L(E)$.
    \end{itemize}
\end{definition}

In regular expressions like in many other algebras there is an assumed precedence of the operators. 

\begin{enumerate}
    \item The $^*$ operator has highest precedence. That is, it applies only to the smallest sequence of symbols to its left. $ab^* \neq (ab)^*$.
    \item Next is concatenation or the "dot" operator. 
    \item Last is the union operator $+$. 
\end{enumerate}

But of course we can manually denote precedence by using parentheses. 

\subsection{Finite automata and Regular expressions}
The languages expressed by regular expressions are precisely the same as those expressed by finite automaton. That is, the regular languages. 

\begin{theorem}
    If $L = L(A)$ for some FA $A$, then there is a regular expression $R$ such that $L = L(R)$.
\end{theorem}

\begin{proof}
    At the heart of this construction is regular expressions denoted as $R^X_{ij}$, that represent strings constructed such that $\hat\delta_A(i, w) = j$ where $w \in \Sigma^*$ and the path between $i \rightarrow j$ (not including $i,j$) does not contain any state $q\in X$.

    \textbf{Basis:} $X = \emptyset$

    When $X = \emptyset$ we are only allowed direct edges between $i$ and $j$, as paths are not allowed any intermediate states. Hence we have two cases. 

    If $i \neq j$ then there are only three cases 
    \begin{enumerate}
        \item There is no symbol $a$ that take us from state $i \rightarrow j$, then $R^\emptyset_{ij} = \emptyset$.
        \item There is exactly one symbol $a$ that label an edge between $i \rightarrow j$, then $R^\emptyset_{ij} = a$.
        \item There are $a_1, \dots, a_k$ symbols that label direct edges between $i \rightarrow j$. Then $R^\emptyset_{ij} = a_1 + \dots + a_k$.
    \end{enumerate}

    If $i = j$ then the legal paths are either of length $0$ (i.e $\epsilon$), or all loops from $i$ to $i$ itself. If the loops on $i$ are labeled by symbols $a_1, \dots, a_k$ then $R^\emptyset_{ii} = \epsilon + a_1 + \dots + a_k$

    \textbf{Induction:} For nonempty $X$. Suppose there is a path between state $i \rightarrow j$ that does not go through state $q \in X$.

    Then there are two cases 
    \begin{enumerate}
        \item The path does not go through state $q$ at all. In this case the label of the path is in the language $R^{X\setminus\set{q}}_{ij}$.
        \item The path goes through state $q$ at least once. Then the path can be split into three pieces: (1) from $i \rightarrow q$, (2) from $q \rightarrow q$ (loops of some length), (3) from $q \rightarrow j$. All of these are in the languages (1) $R^{X\setminus\set{q}}_{iq}$, (2) $R^{X\setminus\set{q}}_{qq}$, (3) $R^{X\setminus\set{q}}_{qj}$.
    \end{enumerate}

    Hence we can form the regular expression 
    $$R^{X}_{ij} = R^{X\setminus\set{q}}_{ij} + R^{X\setminus\set{q}}_{iq} (R^{X\setminus\set{q}}_{qq})^* R^{X\setminus\set{q}}_{qj}$$

    The set of accepted strings for this regular expression is then: 
    $$R = \sum_{f \in F_A} R^{Q}_{q_0f}$$
\end{proof}

\begin{theorem}
    Every language $L = L(R)$ defined by some regular expression $R$ is also defined by a finite automata $A$ s.t $L = L(A)$
\end{theorem}

\begin{proof}
    We shall construct $\epsilon$-NFA's for each case of the regular expression and show by structural induction that we always can come up with such a construction. 

    \textbf{Basis:} Suppose $R$ is some elementary regular expression, that is
    \begin{enumerate}[label=\alph*)]
        \item $R = \set{\epsilon}$
        \item $R = \emptyset$
        \item $R = \set{a}, \quad a\in \Sigma$
    \end{enumerate}

    In Fig.\ref{fig:re-fa-base} we construct the following automata for each case. 
    \begin{figure}[h]
        \centering
        \begin{subfigure}{0.9\textwidth}
            \centering
            \begin{automata}
                \node[state, initial] (q0) {};
                \node[state, right=of q0, accepting] (q1) {};

                \path (q0) edge[] node{$\epsilon$} (q1);
                \node[draw, rectangle, rounded corners, fit=(q0)(q1), inner sep=10pt] {};
            \end{automata}
            \caption{}
        \end{subfigure}
        \begin{subfigure}{0.9\textwidth}
            \centering
            \begin{automata}
                \node[state, initial] (q0) {};
                \node[state, right=of q0, accepting] (q1) {};

                \node[draw, rectangle, rounded corners, fit=(q0)(q1), inner sep=10pt] {};
            \end{automata}
            \caption{}
        \end{subfigure}
        \begin{subfigure}{0.9\textwidth}
            \centering
            \begin{automata}
                \node[state, initial] (q0) {};
                \node[state, right=of q0, accepting] (q1) {};

                \path (q0) edge[] node{$a$} (q1);

                \node[draw, rectangle, rounded corners, fit=(q0)(q1), inner sep=10pt] {};
            \end{automata}
            \caption{}
        \end{subfigure}
        \caption{Construction of $\epsilon$-NFA for the elementary regular expression.}
        \label{fig:re-fa-base}
    \end{figure}

    \textbf{Induction:} We assume the theorem is true for all smaller expressions. Then for the operations 
    \begin{enumerate}[label=\alph*)]
        \item $R + S$
        \item $RS$
        \item $R^*$
    \end{enumerate}

    In Fig.\ref{fig:re-fa-ind} we construct the following automata by the construction of the automata for the sub expressions $R$ and $S$.

    \begin{figure}[h]
    \centering
    \begin{subfigure}{0.9\textwidth}
        \centering
        \begin{automata}
            \node[state, initial] (qs) {};

            \node[state, above right=of qs] (qri) {};
            \node[state, right=of qri] (qre) {};

            \node[state, below right=of qs] (qsi) {};
            \node[state, right=of qsi] (qse) {};
            
            \node[state, below right=of qre, accepting] (qa) {};

            \path (qs) edge[] node{$\epsilon$} (qri);
            \path (qre) edge[] node{$\epsilon$} (qa);

            \path (qs) edge[] node{$\epsilon$} (qsi);
            \path (qse) edge[] node{$\epsilon$} (qa);
            
            \node[draw, rectangle, rounded corners, fit=(qri)(qre), inner sep=10pt] {$R$};
            \node[draw, rectangle, rounded corners, fit=(qsi)(qse), inner sep=10pt] {$S$};
        \end{automata}
        \caption{}
    \end{subfigure}
    \begin{subfigure}{0.9\textwidth}
        \centering
        \begin{automata}
            \node[state, initial] (qs) {};
            
            \node[state, right=of qs] (qri) {};
            \node[state, right=of qri] (qre) {};

            \node[state, right=of qre] (qsi) {};
            \node[state, right=of qsi] (qse) {};
            
            \node[state, right=of qse, accepting] (qa) {};

            \path (qs) edge[] node{$\epsilon$} (qri);
            \path (qre) edge[] node{$\epsilon$} (qsi);
            \path (qse) edge[] node{$\epsilon$} (qa);

            \node[draw, rectangle, rounded corners, fit=(qri)(qre), inner sep=10pt] {$R$};
            \node[draw, rectangle, rounded corners, fit=(qsi)(qse), inner sep=10pt] {$S$};
        \end{automata}
        \caption{}
    \end{subfigure}
    \begin{subfigure}{0.9\textwidth}
        \centering
        \begin{automata}
            \node[state, initial] (qs) {};
            
            \node[state, right=of qs] (qri) {};
            \node[state, right=of qri] (qre) {};
            
            \node[state, right=of qre, accepting] (qa) {};

            \path (qs) edge[] node{$\epsilon$} (qri);
            \path (qre) edge[] node{$\epsilon$} (qa);

            \path (qre) edge[bend right=100] node[above]{$\epsilon$} (qri);

            \path (qs) edge[bend right=40] node{$\epsilon$} (qa);

            

            \node[draw, rectangle, rounded corners, fit=(qri)(qre), inner sep=10pt] {$R$};
        \end{automata}
        \caption{}
    \end{subfigure}
    \caption{Construction of $\epsilon$-NFA for composite regular expressions.}
    \label{fig:re-fa-ind}
\end{figure}
\end{proof}
\FloatBarrier
\subsection{Algebraic Laws for Regular Expressions}

\begin{itemize}
    \item Union commutativity: $L + M = M + L$
    \item Union associativity: $(L + M) + N = L + (M + N)$
    \item Concatenation associative: $(LM)N = L(MN)$
    \item Union identity: $L + \emptyset = \emptyset + L = L$
    \item Concatenation identity: $\epsilon L = L \epsilon = L$
    \item Concatenation annihilator: $\emptyset L = L \emptyset = \emptyset$
    \item Left distributive concatenation over union: $L(M + N) = LM + LN$
    \item Right distributive concatenation over union: $(M + N)L = ML + NL$
    \item Union idempotent: $L + L = L$
    \item Closed closures: $(L^*)^* = L^*$
    \item Closed empty language: $\emptyset^* = \epsilon$
    \item Closed empty string: $\epsilon^* = \epsilon$
    \item Plus closure: $L^+ = LL^*$
    \item Star as Plus: $L^* = L^+ + \epsilon$
    \item Optional: $L? = \epsilon + L$
\end{itemize}

\subsection{Test for RE Algebraic Laws}
% Could elaborate a little, dont feel like it now
To test whether two regular expressions are the same a technique is: Suppose a regular expression $R$ is composed of $L_1,\dots,L_n$, replace each occurrence of $L_i$ by a concrete symbol $w_i$. Do this for both languages and investigate whether they produce the same language. 

